\documentclass[../../main.tex]{subfiles}
\begin{document}

\begin{flushright}
\footnotesize\textit{【自動文字起こし・要確認】}
\end{flushright}

{\bf [解]}
与えられた定積分を
\begin{align}
 F=\int_0^1e^x|x-a|dx \label{eq:1}
\end{align}
とおく．以下$a$の値で場合分けする．

\subparagraph{1. $a\le 0$の時}
\cref{eq:1}より
\begin{align}
  F=-a\int_0^1e^xdx+\int_0^1xe^xdx, \qquad \dfrac{dF}{da}=-\int_0^1e^xdx <0
\end{align}
だから$F$は$a$の単調減少関数である．

\subparagraph{2. $a\ge 1$の時}
\cref{eq:1}より，
\begin{align}
  F=a\int_0^1e^xdx-\int_0^1xe^xdx, \qquad \dfrac{dF}{da}=+\int_0^1e^xdx >0
\end{align}
だから$F$は$a$の単調増加関数である．

\subparagraph{3. $0 \le a\le 1$の時}
以上より$0\le a\le1$の時に$F$は最小値をとるので以下その場合を考える．
$f(x)=e^x(x-a)$とおくと\cref{eq:1}より
\begin{align}
 F=-\int_0^af(x)dx+\int_a^1f(x)dx 
\end{align}
である．$f(x)$の原始関数の一つは$g(x)=e^x(x-a-1)$だから$g'(x)=f(x)$であって，
\begin{align}
\begin{split}
 F 
  &=g(1)+g(0)-2g(a) \\
  &=-ea-a-1+2e^a 
\end{split}
\end{align}
である．よってこれを$a$で微分すると
\begin{align}
  \dfrac{dF}{da}&=-e-1+2e^a 
\end{align}
となるから以下の表を得る．

\begin{table}[htb]
\centering
\caption{$F$の増減表}
\begin{tabular}{c|ccccc}
$a$ & $0$ & & $\log\dfrac{e+1}{2}$ & & $1$ \\
\hline
$F'$ & & $-$ & $0$ & $+$ & \\
\hline
$F$ & & $\searrow$ & & $\nearrow$ &
\end{tabular}
\end{table}
したがって，$a=\log\dfrac{e+1}{2}$で$F$は最小値をとる．

以上の場合分けをまとめて，
\begin{align}
 a=\log\dfrac{e+1}{2}
\end{align}
が求める最小値である．

\newpage

\end{document}
